\section{Термодинамические потенциалы и расчётные зависимости}

Модель IF97 вычисляет свойства через производные
безразмерных термодинамических потенциалов.
Такое представление обеспечивает единую форму записи,
корректное масштабирование переменных
и лучшую численную устойчивость.

В регионах 1, 2 и 5 используется потенциал Гиббса.
В регионе 3 используется потенциал Гельмгольца.
Регион~4 описывается зависимостями насыщения и правилами смешения.

\subsection{Свободная энергия Гиббса для регионов 1, 2 и 5}

Для регионов 1, 2 и 5 стандарт IF97 задаёт безразмерную энергию Гиббса
\begin{equation}
\gamma(\pi, \tau) = \frac{g(p, T)}{R T},
\end{equation}
\begin{gostwhere}
  \item $g$~--- удельная энергия Гиббса;
  \item $R$~--- газовая постоянная воды и водяного пара.
\end{gostwhere}

Безразмерные переменные определяются как
\begin{equation}
\pi = \frac{p}{p^{*}}, \qquad \tau = \frac{T^{*}}{T}.
\end{equation}
Нормировочные константы $p^{*}$ и $T^{*}$ задаются стандартом отдельно для каждого региона.

Для региона~1 функция $\gamma$ задаётся единым полиномиальным разложением
со смещёнными аргументами:
\begin{equation}
\gamma(\pi,\tau) = \sum_{i} n_{i}\,(7{,}1-\pi)^{I_{i}}\,(\tau-1{,}222)^{J_{i}}.
\end{equation}
В этой области используются нормировки
$p^{*}=16{,}53$~МПа и $T^{*}=1386$~К.
Здесь аппроксимация использует $34$ табличных коэффициента.

Для регионов~2 и~5 применяется представление
\begin{equation}
\gamma(\pi,\tau) = \gamma^{o}(\pi,\tau) + \gamma^{r}(\pi,\tau),
\end{equation}
в котором $\gamma^{o}$ описывает идеально-газовую часть,
а $\gamma^{r}$ --- поправку на реальное взаимодействие молекул.
Для региона~2 используются
$p^{*}=1$~МПа и $T^{*}=540$~К,
а для региона~5 ---
$p^{*}=1$~МПа и $T^{*}=1000$~К.
Для региона~2 табличные массивы содержат
$9$ коэффициентов в идеально-газовой части
и $43$ коэффициента в остаточной части.
Для региона~5 используются по $6$ коэффициентов
в идеально-газовой и остаточной частях.

Свойства выражаются через частные производные $\gamma(\pi,\tau)$.
Используются обозначения
$\gamma_{\pi} = \partial \gamma/\partial \pi$,
$\gamma_{\tau} = \partial \gamma/\partial \tau$,
$\gamma_{\pi\pi} = \partial^{2}\gamma/\partial \pi^{2}$,
$\gamma_{\tau\tau} = \partial^{2}\gamma/\partial \tau^{2}$,
$\gamma_{\pi\tau} = \partial^{2}\gamma/\partial \pi \partial \tau$.

Ключевые зависимости, используемые в вычислениях, включают:
\begin{gather}
v = \frac{R T}{p}\,\pi\,\gamma_{\pi}, \\
\rho = \frac{1}{v} = \frac{p}{R T\,\pi\,\gamma_{\pi}}, \\
h = R T\,\tau\,\gamma_{\tau}, \\
s = R \left(\tau\,\gamma_{\tau} - \gamma\right), \\
u = R T \left(\tau\,\gamma_{\tau} - \pi\,\gamma_{\pi}\right), \\
c_{p} = -R\,\tau^{2}\,\gamma_{\tau\tau}, \\
w = \sqrt{R T\,
\frac{\gamma_{\pi}^{2}}
{\frac{\left(\gamma_{\pi} - \tau\gamma_{\pi\tau}\right)^{2}}
{\tau^{2}\gamma_{\tau\tau}}
- \gamma_{\pi\pi}} }.
\end{gather}
Именно эти формулы образуют базовый набор
для прямых расчётов по входной паре $(p, T)$
в регионах~1, 2 и~5~\cite{iapws_if97_2012}.

\subsection{Свободная энергия Гельмгольца для региона 3}

Для региона~3 стандарт использует формулировку через безразмерную энергию Гельмгольца
\begin{equation}
\phi(\delta, \tau) = \frac{f(\rho, T)}{R T},
\end{equation}
\begin{gostwhere}
  \item $f$~--- удельная энергия Гельмгольца.
\end{gostwhere}

Безразмерные переменные определяются как
\begin{equation}
\delta = \frac{\rho}{\rho^{*}}, \qquad \tau = \frac{T^{*}}{T}.
\end{equation}
Для региона~3 фундаментальная зависимость имеет вид
\begin{equation}
\phi(\delta,\tau) = n_{1}\ln\delta + \sum_{i=2}^{40} n_{i}\,\delta^{I_{i}}\,\tau^{J_{i}},
\end{equation}
а нормировочные константы совпадают
с критическими параметрами воды:
$\rho^{*}=\rho_{c}=322$~кг/м\textsuperscript{3} и $T^{*}=T_{c}=647{,}096$~К.

В этой формулировке давление выражается через производную $\phi$ по плотности:
\begin{equation}
p = \rho R T\,\delta\,\phi_{\delta}.
\end{equation}
Это соотношение показывает,
что для региона~3 естественными независимыми переменными
являются $\rho$ и $T$.
При постановке $(p, T)$ требуется сначала определить плотность,
согласованную с заданным давлением.

После определения $\rho$
остальные свойства вычисляются через производные $\phi$:
\begin{gather}
v = \frac{1}{\rho} = \frac{1}{\delta\,\rho^{*}}, \\
u = R T\,\tau\,\phi_{\tau}, \\
s = R \left(\tau\,\phi_{\tau} - \phi\right), \\
h = R T \left(\tau\,\phi_{\tau} + \delta\,\phi_{\delta}\right), \\
c_{p} = R \left[
-\tau^{2}\phi_{\tau\tau}
+ \frac{\left(\delta\phi_{\delta} - \delta\tau\phi_{\delta\tau}\right)^{2}}
{2\delta\phi_{\delta} + \delta^{2}\phi_{\delta\delta}}
\right], \\
w = \sqrt{R T \left[
2\delta\phi_{\delta} + \delta^{2}\phi_{\delta\delta}
- \frac{\left(\delta\phi_{\delta} - \delta\tau\phi_{\delta\tau}\right)^{2}}
{\tau^{2}\phi_{\tau\tau}}
\right]}.
\end{gather}
Формально это такие же табличные аппроксимации,
но в пространстве переменных $(\rho, T)$,
что и делает регион~3 естественным для расчётов
по высокой плотности и околокритическим состояниям~\cite{iapws_if97_2012}.
Важной особенностью является отсутствие
в стандарте аналитических обратных зависимостей для региона~3,
поэтому обратные постановки в этой области
опираются на численные решатели.
